% =============================================================================
%  The TōSh TUI
%  Writing Terminal Interfaces from a Script
% =============================================================================
%
%  A companion to the ToastScript Language Specification. The spec describes the
%  language; this describes one library written in it and for it.
%
%  The preamble deliberately mirrors `toastscript-spec.tex` rather than sharing a
%  file with it. Factoring the two together is worth doing once the second
%  document has settled; doing it now would mean editing a 330 KB document that
%  is under active revision, to save a hundred lines.
% =============================================================================
\documentclass[11pt,a4paper,twoside,openany]{book}

\usepackage[T1]{fontenc}
\usepackage[utf8]{inputenc}
\usepackage{lmodern}
\usepackage{textcomp}
\usepackage[scaled=0.85]{sourcecodepro}
\usepackage{microtype}

\usepackage[top=2.5cm, bottom=2.3cm, inner=3cm, outer=2.5cm, headheight=20pt]{geometry}
\usepackage{parskip}

\usepackage[dvipsnames,svgnames,x11names]{xcolor}
\definecolor{toshblue}{HTML}{2563EB}
\definecolor{toshteal}{HTML}{0891B2}
\definecolor{toshorange}{HTML}{EA580C}
\definecolor{toshgreen}{HTML}{16A34A}
\definecolor{toshpurple}{HTML}{7C3AED}
\definecolor{codebg}{HTML}{F8FAFC}
\definecolor{codeborder}{HTML}{CBD5E1}
\definecolor{codecomment}{HTML}{6B7280}
\definecolor{codestring}{HTML}{059669}
\definecolor{codekeyword}{HTML}{7C3AED}
\definecolor{coverblue}{HTML}{1E3A5F}
\definecolor{covergold}{HTML}{D4A843}
\definecolor{notebg}{HTML}{FEF9C3}
\definecolor{noteborder}{HTML}{CA8A04}
\definecolor{tipbg}{HTML}{DCFCE7}
\definecolor{tipborder}{HTML}{16A34A}
\definecolor{warnbg}{HTML}{FEE2E2}
\definecolor{warnborder}{HTML}{DC2626}
\definecolor{asidebg}{HTML}{EFF6FF}
\definecolor{asideborder}{HTML}{2563EB}

\usepackage{graphicx}
\graphicspath{{assets/}}

\usepackage{booktabs}
\usepackage{longtable}
\usepackage{tabularx}
\usepackage{array}
\newcolumntype{L}[1]{>{\raggedright\arraybackslash}p{#1}}

\usepackage{enumitem}
\setlist{nosep,leftmargin=1.5em}

\usepackage{listings}
\lstdefinelanguage{tosh}{
  keywords={var,const,func,class,record,struct,interface,trait,union,enum,module,
            if,else,for,in,where,while,until,switch,case,default,break,continue,
            return,yield,throw,try,catch,finally,using,new,shy,export,static,prop,
            as,from,ref,out,sealed,fixed,lazy,proud,public,nameof,get,set,true,
            false,null,match,when,subcommand,arg,flag},
  sensitive=true,
  comment=[l]{\#},
  morestring=[b]",
  morestring=[b]',
  literate={|}{{{\color{toshorange}|}}}1
           {=>}{{{\color{toshorange}=>}}}2
           {->}{{{\color{toshorange}->}}}2
           {..}{{{\color{toshorange}..}}}2
           {\$}{{{\color{toshblue}\$}}}1
           {ō}{{\=o}}1
           {│}{{\textbar}}1
           {─}{{-}}1
           {▁}{{\_}}1
           {█}{{\rule{0.5em}{1ex}}}1
           {›}{{>}}1
           {▾}{{v}}1
           {▸}{{>}}1
           {—}{{---}}3
           {–}{{--}}2
           {…}{{...}}3
           {•}{{\textbullet}}1
           {×}{{x}}1
           {°}{{\textdegree}}1,
}
\lstset{
  language=tosh,
  basicstyle=\ttfamily\small,
  keywordstyle=\color{codekeyword}\bfseries,
  commentstyle=\color{codecomment}\itshape,
  stringstyle=\color{codestring},
  backgroundcolor=\color{codebg},
  frame=single,
  rulecolor=\color{codeborder},
  framesep=6pt,
  xleftmargin=8pt,
  xrightmargin=8pt,
  breaklines=true,
  showstringspaces=false,
  tabsize=4,
  aboveskip=\medskipamount,
  belowskip=\medskipamount,
}
\newcommand{\tosh}[1]{\lstinline[language=tosh]{#1}}
\newcommand{\code}[1]{\texttt{\small #1}}

\usepackage[most]{tcolorbox}
\newtcolorbox{notebox}[1][Note]{
  colback=notebg, colframe=noteborder, coltitle=black,
  fonttitle=\bfseries\sffamily\small, title=#1,
  boxrule=0.5pt, arc=2pt, left=6pt, right=6pt, top=4pt, bottom=4pt}
\newtcolorbox{tipbox}[1][Tip]{
  colback=tipbg, colframe=tipborder, coltitle=black,
  fonttitle=\bfseries\sffamily\small, title=#1,
  boxrule=0.5pt, arc=2pt, left=6pt, right=6pt, top=4pt, bottom=4pt}
\newtcolorbox{warnbox}[1][Warning]{
  colback=warnbg, colframe=warnborder, coltitle=black,
  fonttitle=\bfseries\sffamily\small, title=#1,
  boxrule=0.5pt, arc=2pt, left=6pt, right=6pt, top=4pt, bottom=4pt}
\newtcolorbox{aside}[1][]{
  colback=asidebg, colframe=asideborder, coltitle=black,
  fonttitle=\bfseries\sffamily\small, title=#1,
  boxrule=0.5pt, arc=2pt, left=6pt, right=6pt, top=4pt, bottom=4pt}

\usepackage{titlesec}
\titleformat{\chapter}[display]
  {\normalfont\huge\bfseries\sffamily\color{toshblue}}
  {\chaptertitlename~\thechapter}{18pt}{\Huge}
\titleformat{\section}
  {\normalfont\Large\bfseries\sffamily\color{toshblue!90!black}}{\thesection}{1em}{}
\titleformat{\subsection}
  {\normalfont\large\bfseries\sffamily\color{toshteal}}{\thesubsection}{1em}{}
\titleformat{\subsubsection}
  {\normalfont\normalsize\bfseries\sffamily\color{toshpurple}}{\thesubsubsection}{1em}{}

\usepackage{fancyhdr}
\pagestyle{fancy}
\fancyhf{}
\fancyhead[LE]{\sffamily\small\leftmark}
\fancyhead[RO]{\sffamily\small\rightmark}
\fancyfoot[C]{\sffamily\small\thepage}
\renewcommand{\headrulewidth}{0.4pt}
\fancypagestyle{plain}{\fancyhf{}\fancyfoot[C]{\sffamily\small\thepage}\renewcommand{\headrulewidth}{0pt}}

\usepackage[
  colorlinks=true, linkcolor=toshblue, urlcolor=toshteal,
  bookmarks=true, bookmarksnumbered=true, pdfstartview=FitH,
  pdftitle={The T\=oSh TUI}, pdfauthor={The T\=oSh Project},
  pdfsubject={Terminal Interface Framework Reference},
  pdfkeywords={tosh, toastscript, tui, terminal, widgets}
]{hyperref}
\usepackage{bookmark}

\usepackage{caption}
\captionsetup{font=small,labelfont=bf}
\sloppy
\hbadness=10000
\vbadness=10000
\raggedbottom
\usepackage{silence}
\WarningFilter{longtable}{Column widths have changed}
\WarningFilter{latexfont}{Font shape}

\newenvironment{signature}{%
  \par\vspace{4pt}%
  \noindent\textbf{\sffamily\color{toshblue}Signature}\par\nopagebreak
  \begin{quote}\ttfamily\small
}{%
  \end{quote}\vspace{2pt}%
}

\begin{document}

\hypersetup{pageanchor=false}
\begin{titlepage}
\centering
\vspace*{5cm}
{\sffamily\Huge\bfseries\color{coverblue} The T\=oSh TUI\par}
\vspace{0.6cm}
{\sffamily\Large\color{covergold} Writing Terminal Interfaces from a Script\par}
\vspace{2.5cm}
{\sffamily\large A companion to the ToastScript Language Specification\par}
\vfill
{\sffamily\small The T\=oSh Project\par}
\end{titlepage}
\hypersetup{pageanchor=true}

\frontmatter
\tableofcontents

\mainmatter

% =============================================================================
\chapter{What This Is}
% =============================================================================

T\=oSh ships a terminal interface framework, and it is a library rather than a
feature of the shell. The same widgets that draw \code{help browse} and
\code{config browse} are the ones a script builds a window out of, and nothing in
the framework knows it is being driven from a script rather than from C\#.

That is the whole design goal: a T\=oSh script should build a terminal interface
the way a program builds a GTK or Windows Forms window --- construct objects, set
properties, answer events, run it --- and should be able to write the same window
down as a value instead, when keeping the design in its own file is worth more
than building it a line at a time.

\section{Two Ways to Say the Same Thing}

Every screen in this document can be written twice. The two forms produce the
same widget tree and are interchangeable; they are not two frameworks, and one is
not a reduced version of the other.

\subsection{Object style}

Construct widgets, set properties, put them in containers.

\begin{lstlisting}
using Tosh.Tui.Widgets

var name = new TuiField("Name", "world")
var form = new TuiForm(name)

$form.Title = "Greeting"
$form.Submitted = func() { writeline $"Hello, {$name.Text}!" }

tui run $form
\end{lstlisting}

\subsection{Markup style}

Write the tree down as a record. The first key of each node names the widget; the
rest are its properties.

\begin{lstlisting}
tui run {|
    Form  = [ {| Field = "Name", Id = "name", Value = "world" |} ],
    Title = "Greeting",
    OnSubmit = func(form) { writeline $"Hello, {$form.Values.name}!" }
|}
\end{lstlisting}

\begin{notebox}[Why records rather than a file format]
A record literal already nests, already carries a \code{TimeSpan} from
\tosh{1s} and a callable from \tosh{&func}, and already reports its mistakes
through the diagnostics a script author knows. A separate markup language would
be a second parser and a second error reporter, buying none of that. A tree
written as a value can still be moved into a file of its own and brought back
with \tosh{source}, which is the reason to write one this way.
\end{notebox}

\section{Where a Screen Runs}

A full-screen interface needs a terminal. Run a script whose output is a
terminal --- not a pipeline and not a redirect. The framework takes the alternate
screen, hides the cursor, enables mouse reporting, and puts all three back when it
leaves, including when the process is signalled.

\begin{warnbox}
\tosh{tui run} inside a pipeline has nowhere to draw. The screen is not shown and
the command reports that it needs a terminal.
\end{warnbox}

% =============================================================================
\chapter{Running a Screen}
% =============================================================================

\section{\texttt{tui run}}

\begin{signature}
tui run <tree> [--title <text>] [--refresh <duration>]
              [--budget <duration|off>] [--result]
              [--plain [--width <n>] [--height <n>]]
\end{signature}

The tree is a widget built with \tosh{new}, a record literal, or anything
record-shaped. It may be passed as an argument or piped in; the first value shaped
like a tree is the one used, so \code{-{}-title} may sit on either side of it.

\begin{itemize}
  \item \code{-{}-title} draws a titled border around the whole screen. A
        \code{Form} names its own window, so a form needs no title flag.
  \item \code{-{}-refresh} redraws on a timer with nobody touching the keyboard.
        Accepts \tosh{1s}, \code{500ms} or \code{00:00:01}.
  \item \code{-{}-budget} is how long one of your functions may hold the render
        loop. Five seconds by default; \code{-{}-budget off} for a screen that means
        to be slow.
  \item \code{-{}-plain} draws the screen once and answers with its text instead of
        taking over the terminal. \code{-{}-width} and \code{-{}-height} set the
        size; without them it is the terminal's, or 80\,$\times$\,24 when there is
        no terminal at all.
  \item \code{-{}-result} answers with a \code{TuiScreenOutcome} carrying every
        identified widget's value and whether the screen was cancelled.
\end{itemize}

Your bindings and handlers run on the thread that answers the keyboard, which is what
makes them safe to write --- nothing else is touching the widgets while yours does.
The price is that one which never returns is a terminal that never comes back, so each
call carries a deadline. A function that outstays it is stopped, the widget keeps
whatever it last showed, and the reason appears across the bottom of the screen:

\begin{verbatim}
 handler failed - TimeoutException: a screen handler ran for longer than 5s
 and was stopped. Do the slow part in a job and let the screen read the
 result, or pass --budget off.
\end{verbatim}

\noindent Which is the advice: slow work belongs in a job or a channel, with a
\code{Feed} bringing the answer back (page~\pageref{sec:feeds}). Raising the budget
makes a screen wait longer before it tells you something is wrong; it does not make
the screen answer the keyboard while it waits.

A screen containing a \code{Form} emits nothing unless \code{-{}-result} asks: its
handlers have already said what happened, and putting the values on the pipeline
behind the script's own output would be noise.

\section{Asking One Question}

The \code{tui} command has a second job, which is asking a question without
building a screen for it.

\begin{lstlisting}
var name    = tui input "Project name?" --default demo --cli
var go      = tui confirm "Deploy now?" --cli
var chosen  = ls | tui pick --display Name
var path    = tui file --directory
var command = help | tui filter
\end{lstlisting}

\code{-{}-cli} asks at the cursor instead of taking the whole screen, which is what
a shell prompt usually wants.

% =============================================================================
\chapter{Layout}
% =============================================================================

\section{Lengths}

Every widget has a \code{Size}, and a container divides its space according to
what its children ask for. There are four kinds and any of them may be bounded.

\begin{center}
\begin{tabular}{@{}>{\ttfamily\small}l L{9.5cm}@{}}
\toprule
\textbf{\normalfont\sffamily Spelling} & \textbf{\sffamily Meaning} \\
\midrule
auto      & As big as the content. The default. \\
12        & Exactly twelve cells. \\
*         & A share of what is left over. \\
3*        & Three shares of what is left over. \\
25\%      & A quarter of the whole container. \\
1/3       & A third of the whole container, exactly. \\
1/3 24..40 & A third, but never below 24 nor above 40. \\
* ..40    & A share of the leftovers, capped at 40. \\
auto 10.. & As big as the content, but at least 10. \\
\bottomrule
\end{tabular}
\end{center}

\begin{notebox}[A third is not 33\%]
118 columns by 33/100 is 38, and by 1/3 is 39. Percentages are the readable
spelling when the number is round; a ratio is the exact one when it is not.
\end{notebox}

Bounds are what stop a pane being clamped at the call site. Written as a length,
\code{1/3 24..40} stays a third when the terminal is resized; computed with
arithmetic and assigned as a fixed size, it stops being a third the moment either
bound bites.

\section{Rows and Columns}

\begin{lstlisting}
var panes = new TuiStack(TuiOrientation.Horizontal) {|
    Gap   = 1,
    Items = [
        new TuiBorder($list,   "Items")  {| Size = "1/3 24..40" |},
        new TuiBorder($detail, "Detail") {| Size = "*"          |}
    ]
|}
\end{lstlisting}

\begin{lstlisting}
{| Row = [ ... ], Gap = 1 |}
{| Column = [ ... ] |}
\end{lstlisting}

Everything with an answer of its own is settled first, in the order written and
each clamped to its bounds; only then is the remainder shared between the star
children. The rounding error goes to a star child that can absorb it rather than
to one pinned at a bound.

\section{A Status Bar}
\label{sec:statusbar}

There is no status-bar widget, because a row already is one. \code{auto} on the
outer two and \code{*} in the middle puts one group hard left, one hard right, and
one centred between them:

\begin{lstlisting}
{| Row = [
       {| Text = &Name,  Size = "auto" |},
       {| Text = &Clock, Size = "*", Align = "center" |},
       {| Text = &Where, Size = "auto" |}
   ], Size = "1" |}
\end{lstlisting}

Each segment is an ordinary node, so each carries its own style, and each can be
bound --- a clock updates without anything rebuilding the row. Nobody has to know
the terminal's width or measure what they wrote.

\begin{warnbox}[Say \texttt{auto} or share the row]
A child with no \code{Size} shares the row evenly with its siblings, which is right
for panes and wrong for a status bar: the right-hand group ends up floating in the
middle of its third. Write \code{auto} on the segments that should be their own
size.
\end{warnbox}

When the segments do not all fit, every one of them gives up the same proportion of
what it asked for, and what it gives up is marked with an ellipsis. Satisfying them
in written order instead would mean the order decided who got \emph{nothing} --- a
path long enough to fill the row would take all of it and the position indicator at
the other end would silently disappear.

A minimum is how a segment says it matters:

\begin{lstlisting}
{| Text = &Where, Size = "auto 12.." |}
\end{lstlisting}

That one keeps its twelve columns on a narrow terminal and the path gives up the
rest.

\begin{tipbox}
One row of text is a \emph{label}: it is cut and marked rather than wrapped,
whatever \code{Wrap} says, because folding onto a second row that will never be
drawn shows the first few characters and calls that the answer.
\end{tipbox}

\section{Padding and Alignment}

Every widget carries both. They are properties of a widget rather than containers to
wrap one in, which is why Lipgloss layouts --- the styling half of Bubbletea --- read
as short as they do.

\begin{lstlisting}
{| Text = "inset", Padding = 1 |}
{| Text = "inset", Padding = "1 2" |}
{| Text = "inset", Padding = "1 2 3 4" |}
\end{lstlisting}

Written the way CSS writes it: one number for all four sides, two for vertical and
horizontal, four for top, right, bottom and left in that order.

\begin{lstlisting}
{| Text = "over here", Align = "center", VAlign = "middle" |}
\end{lstlisting}

\code{Align} is \code{left}, \code{center}, \code{right} or \code{stretch};
\code{VAlign} is \code{top}, \code{middle}, \code{bottom} or \code{stretch}.
Stretch is the default, because almost everything takes the room it is offered.

Alignment happens \emph{inside} the padding, so a centred widget is centred within
its padding rather than being pushed off by it. A widget knows both rectangles: its
\code{Slot} is what the parent gave it, and its \code{Bounds} is what it actually
paints --- inside the padding and after the alignment.

\begin{notebox}[For anyone writing a widget]
\code{Measure} and \code{Arrange} are not the overridable members; \code{MeasureCore}
and \code{ArrangeCore} are. Padding is added to the measure and taken off the arrange
in the base class, because a widget that forgets to account for it is clipped rather
than inset --- and that is the bug this ships with every time it is left to the widget.
Draw a widget with \code{Paint}, which narrows the surface to its bounds;
\code{Draw} takes a surface already narrowed.
\end{notebox}

% =============================================================================
\chapter{The Widgets}
% =============================================================================

Each widget is given with its markup name, the value that name carries, and the
keys it reads. Every key is optional unless stated.

Four keys are common to all of them: \code{Id}, the name a widget's value is reported
under; \code{Size}, its length; and \code{Padding}, \code{Align} and \code{VAlign},
covered in the previous chapter.

\section{Text and Documents}

\subsection{\texttt{Text} --- wrapped prose}

\begin{lstlisting}
{| Text = "some words", Wrap = true, Bold = true, Foreground = "cyan" |}
\end{lstlisting}

Wraps to the width it is given, keeping the line breaks the author wrote. Set
\code{Wrap = false} to clip instead.

\subsection{\texttt{Lines} --- a document that scrolls itself}

\begin{lstlisting}
{| Lines = ["first", "second", "third"], Scrollbar = true |}
\end{lstlisting}

Lines as written, no wrapping, keeping its own offset. Only the visible rows are
drawn, so a thousand-line page costs the rows on screen rather than a
thousand-row buffer. Arrow keys, \code{Page Up}, \code{Page Down}, \code{Home},
\code{End} and the wheel move it.

\begin{tipbox}
\code{Text} is a paragraph; \code{Lines} is a document. Use \code{Lines} for a
log, a help page, a diff --- anything where the line breaks are the author's and
not the terminal's.
\end{tipbox}

\section{Choosing}

\subsection{\texttt{List} --- one column of items}

\begin{lstlisting}
{| List = $items, Display = "Name", Multi = false,
   OnChange = &Selected, OnSelect = &Opened |}
\end{lstlisting}

\code{Display} names a property to show instead of the item itself, and works on
shell records as well as CLR objects. \code{Multi} turns on ticks, toggled with
\code{Space}. The list keeps its selection in view and scrolls itself.

\subsection{\texttt{Table} --- rows of objects in columns}

\begin{lstlisting}
{| Table = (ls | collect), Borders = "all", Header = true, Scrollbar = true |}
\end{lstlisting}

With no \code{Columns}, the columns are taken from the first row --- and from the
shell's own idea of what is worth showing, so \tosh{ls | tui table} shows the four
columns \tosh{ls} shows at the prompt, with the same alignment and the same cell
text.

\begin{lstlisting}
{| Table = $rows, Columns = [
      "Name",
      {| Header = "Size", Property = "Length", Align = "right", Width = "12" |}
   ] |}
\end{lstlisting}

\code{Borders} is \code{"none"}, \code{"header"} or \code{"all"}; a bare
\tosh{true} means the full grid. None by default, because a table inside a
\code{Box} already has an outline.

Left and right move along the columns when there are more than fit --- whole columns
dropped off the left rather than cells shifted sideways, so every column on screen is
intact and its header still sits above it.

\subsection{\texttt{Tree} --- anything with children}

\begin{lstlisting}
{| Tree = $root, Children = &Kids, Display = &Label, Key = &Path, Root = true |}
\end{lstlisting}

\code{Children} is asked of a node when it opens, not when the tree is built. A
visible closed node is asked whether it has a \emph{first} child, which is what
tells a leaf from an unopened branch.

The guides and markers follow \code{Theme.Tui.TreeStyle}: \code{Clean} draws
indentation and markers alone, \code{Dense} draws connectors between them. A tree that
wants something else says so.

\code{Key} is a node's stable identity, and expansion is remembered against it ---
so a tree rebuilt from fresh objects keeps what the reader had open. \code{Right}
opens and then steps in, \code{Left} closes and then steps out, \code{Enter}
toggles a branch and activates a leaf.

\begin{warnbox}
\code{Key} is not \code{Id}. \code{Id} is the name a widget's value is reported
under; \code{Key} is how a \emph{node} is identified. A tree needs both.
\end{warnbox}

\section{Input}

\subsection{\texttt{Field} --- a caption and an input on one row}

\begin{lstlisting}
{| Field = "Width", Id = "width", Value = "80",
   Placeholder = "columns", Password = false, Multiline = false,
   OnChange = &Typed, OnSubmit = &Entered |}
\end{lstlisting}

The caption is not focusable and the input is, so the keyboard lands on the input
without anything having to say so. The row answers with what its input holds, so
an \code{Id} on the field is the id its text is reported under.

\subsection{\texttt{Combo} --- one of a few}

\begin{lstlisting}
{| Combo = $coolants, Id = "coolant", OnChange = &Picked |}
{| Combo = &Coolants, Display = "Name", Placeholder = "pick one" |}
\end{lstlisting}

A list on the page costs a row per option whether or not anyone is choosing. A
combo box costs one row and opens when it is asked to.

\begin{center}
\begin{tabular}{@{}ll@{}}
\toprule
\textbf{\normalfont\sffamily Key} & \textbf{\normalfont\sffamily What it does} \\
\midrule
Down, Enter, Space, F4 & open the list \\
Up / Down              & move inside it \\
Enter                  & choose, and close \\
Escape                 & close without choosing \\
Left / Right           & step through without opening it at all \\
\bottomrule
\end{tabular}
\end{center}

It is as wide as its widest option, so a column of them lines up and the field does
not change width as the choice does.

\begin{aside}[The same drop-down]
The list is the one a menu bar opens, on the same layer, anchored the same way. That
is why this widget is short and why it behaves the way the menus do: a drop-down
anchored to a field and one anchored to a name on a bar are the same drop-down, and
writing the second one again is what \code{TUI-0023} exists to prevent.
\end{aside}

\subsection{\texttt{Check}}

\begin{lstlisting}
{| Check = "Follow the log", Value = true, OnChange = &Following |}
\end{lstlisting}

Drawn as \code{[x]} or \code{[\ ]} with its label beside it, and a marker in front
when it has the keyboard --- so which one you are on is legible on a terminal with
no colour, the same way a button's is. \code{Enter} and \code{Space} toggle it, and
so does a click.

\code{OnChange} is handed the new state, not the old one, so a handler reads as the
thing it is deciding:

\begin{lstlisting}
func Following(on) { $follow = $on }
\end{lstlisting}

\subsection{\texttt{Group}}

\begin{lstlisting}
{| Group = $settings, Label = "Name", Value = "Current",
   OnToggle = &Flip, OnEdit = &Prompt |}
\end{lstlisting}

A list of named values, drawn as \code{name: value} with a cursor. The arrows,
\code{Home} and \code{End} move it; \code{Space} asks for the row under the cursor
to be toggled, \code{Enter} or \code{e} to edit it, \code{t} to edit it as raw text,
and \code{Esc} closes the group.

Unlike \code{Collection}, this does \emph{not} act on the rows. A collection is a
list of values and ``add this one'' means the same thing everywhere; a group is a
list of settings, and what toggling one means belongs to whoever defined it --- so
the handler is told which row was asked about and decides.

\code{Label} and \code{Value} name the properties read off each row. A row with no
value is drawn as its name alone. This is a list of fields, not a layout: \code{Row},
\code{Column} and \code{Box} are how markup groups widgets on screen.

\subsection{\texttt{Collection}}

\begin{lstlisting}
{| Collection = ["*.tmp", "node_modules"], OnChange = &Save |}
{| Collection = $rules, Display = "Name", Exit = true, OnClose = &Done |}
\end{lstlisting}

A list you can add to, edit and remove from. \code{n} adds an entry, \code{Enter}
or \code{e} edits the one under the cursor, \code{Delete} or \code{r} removes it,
and \code{Esc} closes the editor. While typing, \code{Esc} abandons the line rather
than the editor, and \code{Enter} keeps it.

The widget applies the edit to its own list and then calls \code{OnChange} with the
whole list --- not with the change. A caller almost always wants the result, and one
given only the delta has to keep its own copy in step to find out what it now holds.

\code{Display} names a property to read off each entry for its label. \code{Exit =
true} ends the screen when the editor closes; it is off by default, because an
editor is usually one field among several.

\subsection{\texttt{Choice}}

\begin{lstlisting}
{| Choice = ["Always", "Never", "When asked"], OnChoose = &Decide |}
{| Choice = $themes, Display = "Name", OnChoose = &Apply, OnCancel = &Keep |}
\end{lstlisting}

One choice out of several, drawn as radio rows: the marked option carries
\code{(*)} and the rest \code{( )}, so the reader can see which one \emph{is} the
value rather than only which one the cursor is on.

The arrows, \code{PageUp}, \code{PageDown}, \code{Home} and \code{End} move the
mark. \code{Enter} settles it and calls \code{OnChoose} with the option;
\code{Esc} abandons the choice and calls \code{OnCancel}. \code{OnChange} fires as
the mark moves, before anything is settled.

\code{Display} names a property to read off each option for its label, for a list
of records rather than of strings.

This is not \code{List}. A list is something you move through and activate; this is
a question with one answer, and it keeps that answer's key so redrawing with the
same options leaves the choice where it was.

\subsection{\texttt{Confirm}}

\begin{lstlisting}
{| Confirm = "Discard your changes?", OnAnswer = &Decide |}
{| Confirm = "Delete it?", Title = "Careful",
   ConfirmLabel = "Delete", CancelLabel = "Keep", OnAnswer = &Decide |}
\end{lstlisting}

A question with two answers, drawn in its own frame. \code{OnAnswer} is given
\code{true} for the affirmative answer and \code{false} for the other --- one handler
rather than two, because a dialog that is dismissed has been answered, and a caller
given only an \code{OnConfirm} finds it easy to ignore the cancellation.

Left, right and \code{Tab} move the selection; \code{Enter} takes it. \code{y} and
\code{n} answer outright, as does \code{Esc}, which cancels.

Answering ends the screen, which is the opposite default to \code{Button}: a button
rarely finishes anything and a dialog nearly always does. \code{Exit = false} keeps
one embedded in a larger form. \code{MaxWidth} and \code{MinWidth} bound the frame,
which is sized to the question rather than to the window --- a confirmation that
fills the screen is a screen, and the reader loses the thing they are deciding about.

\subsection{\texttt{Button}}

\begin{lstlisting}
{| Button = "Save", OnPress = &Save |}
{| Button = "Discard", Exit = true, OnPress = &Forget |}
\end{lstlisting}

A button draws a marker and answers \code{Enter} when it has the keyboard, so a
dialog's buttons say nothing about which one is chosen: focus lands on the first
of them and \code{Tab} moves the answer along with it.

\code{Exit = true} ends the screen after the handler has run. Markup never names
the form it is inside, so without it a dialog could raise itself and have no way
to let the reader out; the object style says the same thing with
\code{\$form.Discard()}.

\section{Pictures}
\label{sec:pictures}

\begin{lstlisting}
{| Image = "cover.png", Fit = "letterbox" |}
{| Image = &Chosen, Placeholder = "  nothing here can read this one" |}
\end{lstlisting}

A terminal that can show pixels and a framework that cannot is a framework that has
decided a file browser's third pane is a list of bytes. One cell is two pixels ---
an upper half block with the top pixel's colour in front and the bottom pixel's
behind --- so a pane of $w \times h$ cells is a canvas of $w \times 2h$.

\code{Fit} is \code{letterbox} (the whole picture, right shape, bars where it does
not reach), \code{crop} (fills the room, right shape, edges lost) or \code{stretch}
(fills the room, wrong shape). Letterbox bars are left unpainted rather than
filled, so the picture sits on whatever is behind it.

An image, a frame of a video and a page of a PDF are the same problem, and the
widget knows none of them: it knows pixels. What turns a path into pixels is a hook
the host installs, and the shell's own asks \code{magick}, \code{ffmpeg} and
\code{pdftoppm} --- which are already on the machine of anyone who wants the
feature, so a format nobody here has heard of works the day the reader installs
something that reads it. Results are remembered by path, size and write time, so
arrowing back up a list is instant and a file that changed is read again.

\begin{tipbox}
A path that loads to nothing --- no tool installed, a format none of them know, a
file that is not a picture at all --- leaves the \code{Placeholder} showing. That is
what a preview pane should do for a file it cannot preview, rather than ending the
screen.
\end{tipbox}

\subsection{Real pixels}

Half-blocks are the floor and the part that always works. Where the terminal speaks
the Kitty graphics protocol --- Ghostty, Kitty and WezTerm do --- the same widget
hands over the pixels themselves instead, and the picture is a picture rather than a
mosaic. Nothing in the markup changes; the layout does not either, because the fit is
worked out on the same half-cell canvas whichever way it is drawn.

Which one is used is read from the environment once, at start-up. A reader whose
terminal lies has the last word:

\begin{lstlisting}
TOSH_TUI_GRAPHICS=kitty    # force it on
TOSH_TUI_GRAPHICS=half     # force it off
\end{lstlisting}

\begin{warnbox}[Inside \texttt{tmux} you get half-blocks]
A multiplexer does not know a picture was sent, so it would be drawn over whatever
pane the multiplexer had decided to put there. That is detected and stepped around
rather than left to surprise you.
\end{warnbox}

A placement is kept by the terminal until it is deleted, which is why each one
carries an id: a picture that moved, closed or was covered has to be taken back, and
only something comparing this frame against the last one knows that happened. The
frame diff does both --- cells and placements --- so an unchanged picture is not sent
again, and a screen that ends with one up hands the terminal back clean.

Pictures are scaled down before they are sent. A thousand-pixel photograph in a
thirteen-cell box is four megabytes of base64 down a pipe, to be scaled to two
hundred pixels at the other end; capping it at what a cell could generously be worth
turns that into one.

\subsection{Sixels}

Where the terminal speaks sixels and not Kitty's protocol --- \code{foot},
\code{mlterm}, \code{contour}, \code{mintty} --- the same widget encodes the same
pixels that way instead. Nothing in the markup changes, and Kitty is preferred
where a terminal offers both.

The difference is worth knowing because it reaches the renderer. A Kitty image is
\emph{placed}: it has an id, the terminal keeps it, and removing it is a message. A
sixel is \emph{painted into the screen} like text: there is no id, nothing to
delete, and a picture that goes away is removed by drawing the cells it covered
again. The frame diff does that --- those cells hold exactly what they held last
frame, so nothing else would have noticed they had stopped being right.

\begin{tipbox}
Sixels are smaller on the wire than they look. A preview that costs a megabyte as
Kitty pixels is tens of kilobytes as sixels, because a 216-colour palette and
run-length encoding suit a photograph better than raw bytes do.
\end{tipbox}

\subsection{\texttt{Spinner}}

\begin{lstlisting}
{| Spinner = "Fetching" |}
{| Spinner = "Building", Frames = "line", Spinning = &Busy, Idle = "ok" |}
\end{lstlisting}

\begin{center}
\begin{tabular}{@{}ll@{}}
\toprule
\textbf{\normalfont\sffamily \texttt{Frames}} & \textbf{\normalfont\sffamily What it is} \\
\midrule
\code{dots}     & braille dots --- the default, and what most tools use \\
\code{line}     & a turning line, four frames \\
\code{arc}      & a turning arc \\
\code{circle}   & a quartered circle \\
\code{bounce}   & a dot going round the edge \\
\code{blocks}   & a block growing and shrinking \\
\code{ellipsis} & three dots filling up, slow enough to read \\
\code{ascii}    & the one that works anywhere \\
\bottomrule
\end{tabular}
\end{center}

Each style carries its own pace, because the two are one choice: braille dots read
as motion at eighty milliseconds and as a flicker at forty. A name nobody
recognises spins with dots rather than failing.

The frame is a function of the clock rather than a counter something advances, so
nothing has to tick it and two spinners on one screen cannot disagree about whose
turn it is. What a spinner does need is redraws, and it \emph{asks} for them rather
than starting a thread:

\begin{lstlisting}
tui run {| Spinner = "Working" |}
\end{lstlisting}

\noindent That animates. There is no \code{-{}-refresh}: a screen holding a
spinning spinner reports an interval of its own accord, and stops reporting one the
moment nothing is turning, so a finished screen goes back to costing nothing.

\begin{tipbox}
The text sits at the style's width rather than at the current frame's, so it does
not shuffle left and right as the spinner turns. Several of these characters are
ambiguous-width and draw as two columns on some terminals, which is exactly the
case that would.
\end{tipbox}

\section{Containers}

\subsection{\texttt{Box} --- a titled border}

\begin{lstlisting}
{| Box = [ {| Text = "inside" |} ], Title = "Details" |}
\end{lstlisting}

\subsection{\texttt{Row} and \texttt{Column}}

\begin{lstlisting}
{| Row = [ ... ], Gap = 1 |}
\end{lstlisting}

\subsection{\texttt{Scroll}}

A window onto a child taller than the space available. A list, a table, a tree and
a document all scroll themselves; this is for everything else.

\subsection{\texttt{Grid} --- rows and columns that line up}

\begin{lstlisting}
{| Columns = "auto, *", Grid = [
       [ {| Text = "Name"       |}, {| Text = "tosh"  |} ],
       [ {| Text = "Assemblies" |}, {| Text = "12"    |} ],
       [ {| Text = "Cache"      |}, {| Text = "11 MB" |} ]
   ] |}
\end{lstlisting}

A list of rows, each a list of cells. A row of columns is a stack of stacks until two
rows need the same column widths --- and nested stacks cannot do that, because each
divides its own line and knows nothing about its neighbours. Here the label column
agrees on one width and the values line up under each other.

\code{Columns} and \code{Rows} take the same lengths as \code{Size} does, separated
by commas. Undeclared columns are \code{*} and undeclared rows are \code{auto}, which
is what each axis is usually for: a column divides the width it was given, a row is as
tall as what is in it.

\code{ColumnSpan} and \code{RowSpan} on a cell make it cover several.

\begin{warnbox}[Commas between tracks, not spaces]
A length already spends the space on its own bounds --- \code{auto 12..40} is one
length, not two --- so \code{"auto 12..40, *"} is two tracks and
\code{"auto, 12..40"} is a different two. Getting that backwards produces a layout
that looks plausible and is wrong.
\end{warnbox}

\subsection{\texttt{Dock} --- edges first, the rest to the middle}

\begin{lstlisting}
{| Dock = [
       {| Side = "top",    Text = "Build report" |},
       {| Side = "bottom", Help = ""             |},
       {| Side = "left",   List = &Files         |},
       {| Side = "fill",   Lines = &Output       |}
   ] |}
\end{lstlisting}

The shape almost every full-screen application has. A stack can do it only by nesting
--- a column holding a header, a row holding a sidebar and a body, and a footer ---
which is three containers describing one idea.

Each edge takes what it asks for and the fill takes what is left. With nothing marked
\code{fill}, the last child takes the rest, which is what every other container does
with its last child.

\begin{notebox}[Order decides the corners]
Docked top first, the header runs the full width and the sidebar starts below it:

\begin{verbatim}
TOP                        LEFTTOP
LEFTrest rest rest             rest rest rest
    rest rest rest             rest rest rest
\end{verbatim}

\noindent Written the other way round, the sidebar runs full height and the header
starts beside it. That is not an accident to work around --- it is how you say which
of the two owns the corner.
\end{notebox}

\subsection{\texttt{Tabs} --- several views in one rectangle}

\begin{lstlisting}
{| Tabs = [
       {| Label = "Summary", Lines = &Summary |},
       {| Label = "Log",     Lines = &Log     |},
       {| Label = "Errors",  Text  = "none"   |}
   ] |}
\end{lstlisting}

Each item is an ordinary node with a \code{Label} beside whatever names the widget
--- there is no separate \code{Content} key, because the nesting the language
already has is the node itself.

Left and right move between tabs, and the row is a ring: past the last one is the
first. \code{Home} and \code{End} go to the ends. Not \code{Tab} --- that moves the
keyboard between panes, and a reader who has learned what it means should not find
it meaning two things on one screen.

Only the tab showing is measured, arranged and drawn, so a pane that is expensive to
draw costs nothing while it is behind another. The content is held rather than
rebuilt, so switching back is free and whatever state the child had is still there.

\begin{notebox}[A frame says which tab is showing]
A \code{Box} around tabs, or the screen's own frame, is titled with the tab
showing unless it was told otherwise:

\begin{verbatim}
+ Errors ---------------------+      # follows the selection
|Summary  Log  Errors         |
\end{verbatim}

\noindent \code{-{}-title} wins where it is given, and \code{Title} on the tabs
themselves overrides the label without needing a box. A screen whose root is a set of
tabs gets a frame even with no \code{-{}-title}, because a row of labels with nothing
naming the pane beneath them reads as two unrelated things.
\end{notebox}

\subsection{\texttt{Menu} and \texttt{MenuBar}}

A screen can offer a command two ways without a menu. It can register a key, which
is terse and costs no room but is invisible until someone reads the footer --- and
there are not enough memorable chords for five menus of a dozen. Or it can put a
button on the page, which is discoverable and costs a row of the document. A menu
is the third answer: one row, whether there are two menus or nine, because the
group opens only when it is asked for.

\begin{lstlisting}
{| Id = "menus", MenuBar = [
       {| Menu = "File", Items = [
              {| Item = "Open", Key = "Ctrl+O", Do = &Open |},
              {| Separator = true |},
              {| Item = "Quit", Key = "q" |}
          ] |},
       {| Menu = "Edit", Items = [
              {| Item = "Copy", Key = "Ctrl+C", Do = &Copy |},
              {| Item = "Replace", Enabled = false |}
          ] |}
   ], Size = "1" |}
\end{lstlisting}

\begin{center}
\begin{tabular}{@{}ll@{}}
\toprule
\textbf{\normalfont\sffamily Key} & \textbf{\normalfont\sffamily What it does} \\
\midrule
F10 (or a click)   & opens the bar --- see below \\
Left / Right       & walk the bar, and take the open menu with them \\
a letter           & opens the menu it names \\
Up / Down          & move inside the open menu, over separators and disabled commands \\
Enter              & choose, which also takes the menu down \\
Escape             & close, and hand the keyboard back \\
\bottomrule
\end{tabular}
\end{center}

\begin{warnbox}[The key beside a command is shown and not registered]
It belongs to the screen, so register it in \code{Keys}. A key registered on the
menu would answer only while the menu was open, which is exactly backwards ---
\code{Ctrl+O} exists so that you do not have to open the menu.
\end{warnbox}

A bar does not take the keyboard until something is down. It sits at the top of the
tree, so a bar that took focus like any other widget would take it \emph{first}, and
a letter typed at a screen that had just opened would open a menu instead of being
typed. Reach it the way every editor does --- a key aimed at it:

\begin{lstlisting}
Keys = [ {| Key = "f10", Does = "menu", Press = "menus" |} ]
\end{lstlisting}

The popup is a layer, not part of the bar: a widget cannot paint outside what it was
given, which is the whole reason layers exist. A screen with a bar in it gets one
without asking, so the bar is written where it belongs --- one row at the top of a
column --- rather than having the whole screen wrapped around it.

\subsection{\texttt{Form} --- a window with a way in and a way out}

\begin{lstlisting}
{| Form = [ ... ], Title = "Edit", OnSubmit = &Apply, OnCancel = &Bail |}
\end{lstlisting}

Covered in full in Chapter~\ref{ch:events}.

\subsection{\texttt{Overlay} --- one thing on top of another}

\begin{lstlisting}
{| Overlay = [ $page, $dialog ], Margin = 4 |}
\end{lstlisting}

The first child is the page; the second sits on top of it. While something is on
top it is the only thing the keyboard and the mouse can reach, and the content
beneath keeps drawing and stops responding --- which is what \emph{modal} means.

\section{Charts}

\subsection{\texttt{Gauge} --- how far along}

\begin{lstlisting}
{| Gauge = null, Value = 42, Min = 0, Max = 100, Label = "" |}
\end{lstlisting}

The label is drawn over the bar, so a gauge is one row however long its caption
is. \code{Label} unset shows the percentage; an empty label shows nothing.

\subsection{\texttt{Bars} --- several values compared}

\begin{lstlisting}
{| Bars = [ {| Label = "disk", Value = 40 |}, {| Label = "swap", Value = 5 |} ],
   Scale = 100, Values = true |}
\end{lstlisting}

Where a \code{Gauge} says how far along one thing is, this says how several compare.
The labels share a column so the bars start at the same place, and every bar is
measured against the same maximum --- otherwise comparing them says nothing.

A record whose fields are numbers works as well as a list of rows, which reads better
for a fixed set:

\begin{lstlisting}
{| Bars = {| disk = 40, swap = 5, tmp = 12 |} |}
\end{lstlisting}

\subsection{\texttt{Spark} --- a series in one row}

\begin{lstlisting}
{| Spark = $samples, Scale = 100, Capacity = 240, Pad = true |}
\end{lstlisting}

The chart is a window onto the \emph{tail} of the series, so a monitor pushes
samples forever and the widget draws the last however-many fit.

\begin{notebox}[\texttt{Scale} is the one to know about]
\code{Scale} is the lowest the top of the chart will go. Set it to 100 for a
percentage. Without it, a quiet series of 1--2--1--2 scales to its own maximum and
draws as a mountain range, which is precisely the reading a monitor must not give.
\end{notebox}

\subsection{\texttt{Chart} --- a series against a scale you can read}

\begin{lstlisting}
{| Chart = $samples, Scale = 100, Marks = "braille" |}
\end{lstlisting}

The other three have no axis on purpose. This one does, and the axis is the reason to
reach for it: without numbers a chart is a decoration. Ticks are chosen to be numbers
somebody would have written --- 0, 25, 50, 75, 100 rather than 0, 33.3, 66.6 --- and
the labels take their precision from the step, so the column lines up.

\begin{center}
\begin{tabular}{@{}lll@{}}
\toprule
\code{Marks} & Draws & Reach for it when \\
\midrule
\code{blocks} & An area, eighths of a cell & Always, unless you need better \\
\code{braille} & A line, eight points per cell & The shape between samples matters \\
\code{ascii} & One \code{\#} per point & The terminal is known to be poor \\
\bottomrule
\end{tabular}
\end{center}

\noindent \code{Stretch} is on by default and spreads the whole series across the
plot: eight build times in a forty-column pane are five columns each. Turn it off for
a live meter and the plot becomes a window onto the newest samples, one per column,
the way \code{Spark} works. Where there are more samples than columns, each column
takes the \emph{largest} value in its bucket --- a mean is defensible and wrong here,
because a chart that averages away the spike hides the incident the reader opened it
to find.

\begin{notebox}[Which of the four?]
\code{Gauge} is one proportion, now. \code{Bars} compares a handful of labelled
values. \code{Spark} is a shape in one row, for a corner of a status line.
\code{Chart} is the one that takes a pane and answers ``how much, exactly''.
\end{notebox}

\noindent \code{examples/chart.tosh} draws one series all three ways.

% =============================================================================
\chapter{Values and Bindings}
% =============================================================================

\section{Identifying a Widget}

A widget with an \code{Id} reports its value. A form collects them, and
\code{-{}-result} hands them back.

\begin{lstlisting}
$form.Values.Width        # from inside a handler
$answers.Values["Width"]  # from a --result outcome
\end{lstlisting}

\section{Pull Bindings}

A property set to a function is a \emph{question} rather than a value. It is asked
again before every redraw, and a keystroke or a refresh tick is what causes one.

\begin{lstlisting}
var preview = new TuiTextWidget()
$preview.TextSource = &Summary
\end{lstlisting}

\begin{lstlisting}
{| Text = &Summary |}
\end{lstlisting}

In markup the binding is written by putting a function reference where a value
would go; the framework notices that a callable is not content.

A binding is handed the form's current values as a record, so a pane that follows
another widget is a function of the screen's state rather than something that has
to go looking for it:

\begin{lstlisting}
func Summary(values) {
    return $"{$values.Width} x {$values.Height}"
}
\end{lstlisting}

The bindable properties are \code{TextSource} on text, \code{ValueSource} on an
input, \code{ItemsSource} on a list, \code{RowsSource} on a table,
\code{LinesSource} on a document, \code{ValuesSource} on a sparkline,
\code{AmountSource} on a gauge, and \code{TitleSource} on a border.

\begin{tipbox}
There is no update code in a screen built this way, because there is nothing to
keep in step. A live monitor is two pull bindings and \code{-{}-refresh 1s}.
\end{tipbox}

A binding is asked for \emph{everything} it produces, not for the last of it --- and
that is not a convenience, it is the only thing a caller can safely do.

A collection written as an expression is a sequence, and the engine hands one over
two ways. A function returning a bare variable yields the collection's
\emph{elements}, because a variable reference is expanded where it is produced. A
function returning a literal yields the collection \emph{whole}, marked as a
sequence for whatever consumes it to spread. Both are correct, both are
long-standing, and no script can tell them apart --- \code{Items | collect} counts
three either way.

A binding reads the raw stream, where that marking is not visible, so it takes the
lot. Keeping only the last value meant keeping only the last \emph{row}, and every
binding that carries more than one thing --- \code{Lines}, \code{List},
\code{Table}, \code{Spark}, \code{Bars} --- silently showed its final element.

\begin{warnbox}[A literal in a test proves less than it looks]
\code{func Items() => ["a", "b"]} and \code{func Items() => \$rows} exercise
different halves of this. A test written with the literal --- which is what one
naturally writes --- cannot catch a fault in the expanded half, and did not
(\code{TUI-0027}).
\end{warnbox}

\section{Feeds}
\label{sec:feeds}

A refresh interval is polling, and polling is right for a \emph{sample}: CPU usage
is a measurement, not an event. It is the wrong shape for everything else a shell
shows. A log being tailed, a build emitting lines, a download reporting progress:
those arrive when they arrive, and polling them means either a short interval that
burns work re-reading nothing or a long one that lags the event.

A \code{Feed} says where a screen's changes come from:

\begin{lstlisting}
var log = channel

async {
    for step in $steps { channel-send $log $step }
    channel-close $log
} | ignore

tui run {|
    Lines = &Done,
    Feed  = {| Source = $log, Do = &Arrived, OnEnd = &Finished |}
|}
\end{lstlisting}

\begin{lstlisting}
$pane.Feeds = new TuiFeeds()
$pane.Feeds.From($log, &Arrived, &Finished) | ignore
\end{lstlisting}

A source is a channel, an async sequence, a task, or any sequence at all. It is
read on a task of its own, and each value is handed to \code{Do} \emph{on the render
loop's own thread}, between frames. So a feed handler is written exactly like one
answering a keystroke: no lock, no thread to think about, nothing half-drawn, and
a redraw when it returns.

\code{Feed} hangs on the node it belongs to, the way \code{Keys} does, so a screen
with several panes fed by several sources says so on each of them.

\begin{aside}[What this costs when you are not using it]
Nothing. A screen with no feed and no interval blocks on the keyboard exactly as it
always did. Only a screen that has one waits on both, and redraws coalesce --- a
thousand values arriving between two frames produce one frame, because the signal
is a flag rather than a count.
\end{aside}

Interval and feed compose on one screen, and should: the clock in a corner is a
sample and the log beneath it is not.

% =============================================================================
\chapter{Events and Forms}
\label{ch:events}
% =============================================================================

\section{Handlers}

A script function assigned where a handler is wanted becomes one. This is a
language feature rather than a TUI one: a function knows which engine defined it,
so it can be handed to any .NET type that expects a delegate.

\begin{lstlisting}
$field.Changed = func(text) { writeline $"now: {$text}" }
$field.Changed = &Remember
$button.Pressed = func() { $saved = true }
\end{lstlisting}

A handler that ignores its arguments need not declare them. A handler that must
produce a value has what it produced converted rather than cast, so returning
\tosh{1} where a \code{double} is wanted is not a mistake.

\section{Forms}

A form is a window with a way in and a way out. It owns \code{Enter} and
\code{Escape}, runs its handler while the widgets still exist, and reports how it
ended so a script does not have to infer it.

\begin{lstlisting}
var width = new TuiField("Width", "2")
var form  = new TuiForm(width) {| Title = "New block" |}

$form.Submitted = func() { $blockWidth = Whole($width.Text, 2) }

tui run $form

if ($form.WasCancelled) { throw "Cancelled." }
\end{lstlisting}

Handlers are handed the form, the way a Windows Forms event is handed its sender.
A handler holding the widgets in variables ignores it; one built from markup does
not have them and reads \code{\$form.Values.Width} instead.

\subsection{Changing its mind}

\code{KeepOpen} is \code{FormClosing} with \code{e.Cancel}. The result is set
before a handler runs, so a handler can read how the form ended and decide that it
has not:

\begin{lstlisting}
$form.Cancelled = func() {
    $overlay.Modal = $confirm
    $form.KeepOpen()
}
\end{lstlisting}

\code{Discard} is the other half of it. The button that answers the question must
not ask it again --- \code{Cancel} would, and the same dialog would go straight
back up:

\begin{lstlisting}
$discard.Pressed = func() { $form.Discard() }
\end{lstlisting}

\subsection{Keys travel outwards}

\code{Enter} and \code{Escape} reach a form only after the focused widget and
everything between have declined them --- so a multi-line field keeps its
newlines, and a list with an activation handler keeps its \code{Enter}.

% =============================================================================
\chapter{Input and Focus}
% =============================================================================

\section{Routing}

The focused widget sees an event first. If it declines, the event travels outwards
to its ancestors, and finally to the screen. A text field consumes printable
characters, so a screen-level shortcut never sees the \code{q} someone is typing
into a search box --- without the screen having to know a text field exists.

\begin{aside}[Why this is structural]
Both browsers once checked their quit key before they checked where focus was, so
\code{q} closed the window mid-search. That is what happens when a shortcut is a
line of code in the middle of a method: whether it fires depends on where the line
sits rather than on what it means.
\end{aside}

\section{Focus Scopes}

An overlay showing something modal answers with that alone when asked what the
keyboard may reach. \code{Tab} cannot step behind a dialog, and a click aimed past
one is swallowed rather than landing on a widget the reader cannot currently use.

When a handler puts a dialog up, the keyboard moves into it; when the dialog
closes, it moves back. Only a widget that has gone out of reach is re-seated ---
nothing focused is a state a screen can ask for.

\section{What the Terminal Can Show}

A screen is written once and read on whatever the reader has. The framework asks the
environment what that is, once per session, and reduces on the way out --- so nothing
you write has to know about it.

\subsection{Colour}

Colour degrades rather than disappearing. A named colour is already one of the
sixteen every terminal has, so most of a theme survives untouched; a hex colour
becomes a palette index at 256 colours and the nearest of sixteen below that.

With no colour at all, a \emph{background} becomes reverse video. That is the
difference between degrading and dropping: a selected row is drawn with a background,
and a reader who cannot see which row is selected has lost something that matters.

\begin{center}
\begin{tabular}{@{}ll@{}}
\toprule
Variable & Effect \\
\midrule
\code{NO\_COLOR} & Any non-empty value turns colour off. \\
\code{FORCE\_COLOR} & \code{0}--\code{3}: none, 16, 256, truecolor. \\
\code{COLORTERM} & \code{truecolor} or \code{24bit} means 24-bit colour. \\
\code{TOSH\_TUI\_COLOR} & \code{none}, \code{16}, \code{256}, \code{truecolor} --- wins over all of them. \\
\bottomrule
\end{tabular}
\end{center}

\noindent Otherwise \code{TERM} decides: \code{256color} in the name means 256,
\code{dumb} or unset means none, anything else means sixteen.

\subsection{Characters}

Borders, tree guides, meters and spinners are drawn out of characters above U+2000.
Where the encoding cannot carry them --- \code{LANG=C}, a serial console, \code{TERM=dumb}
--- they fold to ASCII on the way out:

\begin{verbatim}
+ colour ----------------------+      # was a rounded border
|Disk    72%                   |
|#########:::::::::::::::::::::|      # was block characters
\end{verbatim}

\noindent Only what the framework draws is folded. A filename with a CJK character in
it is your data and is left exactly as it is. \code{TOSH\_TUI\_ASCII=1} forces it.

\begin{warnbox}[A fold never changes a width]
Every replacement is one column, because the layout has already reserved the columns.
This is why \code{…} folds to \code{.} and not to \code{...}, and why a two-column
character is never folded at all --- a fold that changed a width would leave the row
short and every border after it out of line, which is the fault this exists to prevent.
\end{warnbox}

\section{Shortcuts}

Screen-level keys belong in a table, asked once after the focused widget has
declined:

\begin{lstlisting}
var keys = new TuiShortcuts()
$keys.On('/', "/", "search", func() { $focus = "search" })
$keys.Exit(ConsoleKey.Q, "q", "quit")
\end{lstlisting}

Each binding carries the label and description a footer would print, and
\code{Describe()} builds that footer --- so the help line cannot claim a key that
does not work.

\begin{warnbox}[Match punctuation on the character]
.NET reports \code{ConsoleKey.Divide} for \code{/} on a real terminal, and
\code{ConsoleKey.None} for \code{[} and \code{]}. A shortcut written against
\code{Oem2} never fires. Register the character the reader actually pressed.
\end{warnbox}

A chord is written the way a reader would say it: \code{"q"}, \code{"ctrl+s"},
\code{"f5"}, \code{"/"}. \code{"space"}, \code{"esc"}, \code{"up"}, \code{"pgdn"}
and their like are accepted alongside the enum's own names --- and \code{"space"} is
the only way to say it, because a chord is trimmed before it is read and a literal
\code{" "} does not survive that.

\begin{warnbox}[Shift is not a chord]
Shift is how a capital arrives, so a table that treated it as a modifier could not
register \code{r} and \code{R} as two different bindings. Matching ignores it, which
means \code{"shift+tab"} also fires on a plain \code{Tab}. When the press decides
what the key does, take it as an argument:

\begin{lstlisting}
$keys.On(ConsoleKey.Tab, "Tab", "switch panes", func(key) {
    Cycle(reverse: $key.Modifiers.HasFlag(ConsoleModifiers.Shift))
    return TuiScreenResult.Continue
})
\end{lstlisting}

\noindent One entry, one line in the footer, and both directions work.
\end{warnbox}

\subsection{A key somebody else answers}

Not every key on a screen is the screen's. A text editor answers its own
\code{Enter}, and that is right --- the widget that owns the state owns the key.
But the footer still has to say so, and a screen that re-registered the editor's keys
in its own table to be able to describe them would have two dispatchers for one key,
which is worse than the footer it was trying to fix.

\code{Note} describes without dispatching:

\begin{lstlisting}
$keys.When(func() { return $editing }, func() {
    $keys.Note("Enter", "stage")
    $keys.Note("Esc", "cancel")
})
\end{lstlisting}

\noindent A note never fires. What it buys is that the condition under which the key
is \emph{described} is the same expression the dispatcher switches on, so the two
cannot disagree about when --- which is the drift a second hand-written string
eventually produces.

\subsection{A key that aims at the screen}

Running a function covers every command that changes state --- save, open, quit,
refresh --- and none that change where the keyboard \emph{is}, which is the other
half of what keybindings are for. \code{Ctrl+P} opens a palette \emph{and puts the
caret in it}. \code{/} focuses the search box. \code{F2} moves to the second pane.

A key says what it wants by id, the same way everything else on a screen is
addressed:

\begin{lstlisting}
Keys = [
    {| Key = "/",      Does = "search",  Focus = "query"  |},
    {| Key = "escape", Does = "back",    Back  = true     |},
    {| Key = "f5",     Does = "refresh", Press = "reload" |},
    {| Key = "q",      Does = "quit",    Exit  = true     |}
]
\end{lstlisting}

\begin{lstlisting}
$keys.Focus("/", "/", "search", "query") | ignore
$keys.Press("f5", "F5", "refresh", "reload") | ignore
$keys.Back("escape", "Esc", "back") | ignore
\end{lstlisting}

\begin{center}
\begin{tabular}{@{}ll@{}}
\toprule
\textbf{\normalfont\sffamily Key} & \textbf{\normalfont\sffamily What it does} \\
\midrule
\code{Do}    & runs a function \\
\code{Focus} & moves the keyboard to a widget, by id \\
\code{Press} & does what Enter on that widget would do, without moving the keyboard \\
\code{Back}  & puts the keyboard where it was before the last \code{Focus} \\
\code{Exit}  & ends the screen, and composes with \code{Do} \\
\bottomrule
\end{tabular}
\end{center}

\code{Press} deliberately leaves the keyboard alone: \code{F5} should not take the
caret out of whatever the reader was typing in. \code{Back} is what closing a
palette, a menu or a search box wants.

An id written on a labelled field names the row, because that is the node you
wrote --- \code{\{| Field = "Find", Id = "query" |\}} --- and aiming looks inside it
for the input, which is the reading bindings already take. Aiming at something
hidden, missing, or behind a modal does nothing at all; a binding that does not
currently apply should be inert, not fatal.

\begin{warnbox}[A printable key still belongs to whoever is typing]
\code{/} reaches the table only after the focused widget declines it. With a list
focused it opens the search box; with a text field focused it types a slash. That
is the routing working, and it is why a screen key can be a bare letter at all.
\end{warnbox}

\subsection{\texttt{Help} --- a footer nobody maintains}

A footer used to be a string kept beside the \code{switch} that implemented the
keys, and it drifted the way that always does: the help browser offered \code{/},
\code{[} and \code{]} for a while after they had stopped working.

\code{Help} is a projection instead. It shows the keys that would answer
\emph{right now}:

\begin{lstlisting}
{| Help = "", Size = "auto" |}
{| Help = "", Full = &Everything |}
\end{lstlisting}

\begin{center}
\begin{tabular}{@{}ll@{}}
\toprule
\textbf{\normalfont\sffamily It shows} & \textbf{\normalfont\sffamily Because} \\
\midrule
the keys on the focused widget first & that is the one that would answer \\
then each ancestor's, then the screen's & which is the order a key travels \\
a chord bound twice, once             & two lines for one key describes nobody's screen \\
nothing a condition excludes          & a footer must not offer a key that does nothing \\
\bottomrule
\end{tabular}
\end{center}

A condition is the other half of a binding. \code{When} on a key says when it
applies, and a key that does not apply is neither described \emph{nor fired}:

\begin{lstlisting}
Keys = [
    {| Key = "e",      Does = "edit", Do = &Edit, When = &NotEditing |},
    {| Key = "ctrl+s", Does = "save", Do = &Save, When = &Editing    |}
]
\end{lstlisting}

\begin{lstlisting}
$keys.When(func() => not $editing, func() {
    $keys.On("e", "e", "edit", &Edit) | ignore
}) | ignore
\end{lstlisting}

\noindent In code the condition wraps a group rather than being an argument to
each key, because keys that apply together are written together and repeating the
predicate six times is how the seventh comes to disagree. Conditions nest: a group
inside a group applies only where both hold.

\begin{tipbox}
A condition that throws is taken as "no". A footer is not worth ending a screen
over, and a binding whose condition cannot be evaluated is one nobody should be
told about either.
\end{tipbox}

\subsection{Starting something in the background}

A key's handler runs on the render loop's thread, so it should return promptly. Work
that takes time belongs in a job or an \code{async} block whose results come back
through a \code{Feed} (\S\ref{sec:feeds}) --- which is also what keeps the screen
answering keys while it runs.

% =============================================================================
\chapter{Styling}
% =============================================================================

\section{A Style}

\begin{lstlisting}
using Tosh.Tui.Rendering

var warn = new TuiStyle("red", null, TuiTextAttributes.Bold)
\end{lstlisting}

In markup a style is written as keys on the node itself:

\begin{lstlisting}
{| Text = "careful", Foreground = "red", Bold = true, Dim = false |}
\end{lstlisting}

The keys are \code{Foreground}, \code{Background}, \code{Bold}, \code{Dim},
\code{Italic} and \code{Underline}.

\section{Styled Runs}

A row is rarely one style all the way across. A \code{TuiSpanLine} is a row of
\code{TuiSpan} --- text and a style --- and it is what a list or a document draws
when one style per row is not enough.

\begin{lstlisting}
var line = new TuiSpanLine()
$line.Add("› ", $gutter)
$line.Add("grep", $bold)
$line.Add("  [BuiltIn]", $dim)
\end{lstlisting}

\code{SpanSelector} on a list and \code{LineSelector} on a tree hand a caller the
row and take one of these back. Given either, the widget draws no marker of its
own: a row that says how to draw itself draws itself entirely.

\section{Following the Theme}

A screen that should look like the rest of the shell reads the user's
configuration:

\begin{lstlisting}
using Tosh.Tui.Rendering

var theme = $tosh.Config.Theme.Tui
$border.Style = $theme.Border.ToStyle()
\end{lstlisting}

% =============================================================================
\chapter{A Worked Example}
% =============================================================================

A live system monitor: a list of processes on the left, meters and a summary on
the right, redrawn once a second with no update code.

\begin{lstlisting}
using Tosh.Tui.Widgets = Tui
using Tosh.Tui.Rendering

var cpuGauge = new Tui.TuiGauge()     {| Label = "", Size = "2*" |}
var cpuSpark = new Tui.TuiSparkline() {| MinimumScale = 100, Capacity = 240, Size = "3*" |}

func Heat(value: double, ratio: double) -> TuiStyle {
    if ($ratio >= 0.75) { return new TuiStyle("red") }
    if ($ratio >= 0.40) { return new TuiStyle("yellow") }
    return new TuiStyle("cyan")
}

$cpuSpark.StyleSelector = &Heat

## Asked once per redraw, which is the cadence a chart wants — so the sampling
## happens here and the meters are fed from it.
func Summary() {
    var load = $cpu.Sample()

    $cpuGauge.Amount = $load
    $cpuSpark.Push($load)

    return $"Load {$load.ToString("F1")} %"
}

var screen = new Tui.TuiStack(Tui.TuiOrientation.Horizontal) {|
    Items = [
        new Tui.TuiBorder(
            new Tui.TuiList() {| DisplayProperty = "Name", ItemsSource = &TopProcesses |},
            "Processes") {| Size = "2*" |},

        new Tui.TuiBorder(
            new Tui.TuiStack(Tui.TuiOrientation.Vertical) {|
                Items = [
                    $cpuGauge,
                    $cpuSpark,
                    new Tui.TuiTextWidget() {| TextSource = &Summary, Size = "*" |}
                ]
            |},
            "System") {| Size = "3*" |}
    ]
|}

tui run $screen --title "System Monitor" --refresh 1s
\end{lstlisting}

The full version ships as \code{examples/system-monitor.tosh}. Its companion,
\code{examples/reactor-block.tosh}, is a form with a live preview, and
\code{examples/reactor-block-markup.tosh} is the same screen written as markup ---
the two are identical outside the form, so the difference between the styles is
all that separates them.

% =============================================================================
\appendix
\chapter{Markup Reference}
% =============================================================================

Every markup name, what its own key carries, and the keys it reads. \code{Id} and
\code{Size} are read by all of them.

\begin{center}
\small
\begin{longtable}{@{}>{\ttfamily}l >{\ttfamily\footnotesize}L{3.6cm} >{\ttfamily\footnotesize}L{6.4cm}@{}}
\toprule
\textbf{\normalfont\sffamily Name} & \textbf{\normalfont\sffamily Carries} & \textbf{\normalfont\sffamily Keys} \\
\midrule
\endfirsthead
\toprule
\textbf{\normalfont\sffamily Name} & \textbf{\normalfont\sffamily Carries} & \textbf{\normalfont\sffamily Keys} \\
\midrule
\endhead
Text    & the text            & Wrap, + style \\
Lines   & the lines           & Scrollbar, + style \\
List    & the items           & Display, Multi, OnChange, OnSelect, + style \\
Table   & the rows            & Columns, Borders, Header, Scrollbar, OnChange, OnSelect \\
Tree    & the root            & Children, Display, Key, Root, Scrollbar, OnChange, OnSelect \\
Field   & the caption         & Value, Placeholder, Password, Multiline, OnChange, OnSubmit \\
Check   & the label           & Value, OnChange, + style \\
Combo   & the options         & Display, Placeholder, OnChange, + style \\
Button  & the label           & OnPress, Exit, + style \\
Gauge   & ---                 & Value, Min, Max, Label, + style \\
Spark   & the series          & Scale, Capacity, Pad, + style \\
Bars    & the values          & Scale, Values, + style \\
Image   & the path            & Fit, Placeholder, + style \\
Spinner & the text            & Frames, Spinning, Idle, + style \\
Row     & the children        & Gap \\
Column  & the children        & Gap \\
Box     & the children        & Title, + style \\
Help    & ---                 & Full, Separator, + style \\
Scroll  & the child           & --- \\
Menu    & the name            & Items, + style \\
MenuBar & the menus           & Gap, + style \\
Form    & the children        & Title, OnSubmit, OnCancel \\
Overlay & page, then modal    & Margin \\
\bottomrule
\end{longtable}
\end{center}

\vspace{0.5em}
\noindent\textbf{\sffamily On every node:} \code{Id}, \code{Size}, \code{Padding},
\code{Align}, \code{Valign}, \code{Visible}, \code{When}, \code{Keys}, \code{Feed}.

\vspace{1em}
\noindent\textbf{\sffamily Style keys:} \code{Foreground}, \code{Background},
\code{Bold}, \code{Dim}, \code{Italic}, \code{Underline}.

\vspace{0.5em}
\noindent\textbf{\sffamily Common to every name:} \code{Id}, \code{Size},
\code{Padding}, \code{Align}, \code{VAlign} --- see Chapter~3.

\begin{notebox}[Adding a widget]
The registry is open. \code{TuiWidgetRegistry.Register} adds a name, and a widget
reachable only from the object style is considered a gap --- the test suite
compares the registry against the widget types and fails when one has no name.
\end{notebox}

\end{document}
